\documentclass{article}
\usepackage{mathtools} 
\usepackage{fontspec}
\usepackage[UTF8]{ctex}
\usepackage{amsthm}
\usepackage{mdframed}
\usepackage{xcolor}
\usepackage{amssymb}
\usepackage{amsmath}


% 定义新的带灰色背景的说明环境 zremark
\newmdtheoremenv[
  backgroundcolor=gray!10,
  % 边框与背景一致，边框线会消失
  linecolor=gray!10
]{zremark}{说明}

% 通用矩阵命令: \flexmatrix{矩阵名}{元素符号}{行数}{列数}
\newcommand{\flexmatrix}[4]{
  \[
  #1 = \begin{pmatrix}
    #2_{11}     & #2_{12}     & \cdots & #2_{1#4}   \\
    #2_{21}     & #2_{22}     & \cdots & #2_{2#4}   \\
    \vdots      & \vdots      & \ddots & \vdots     \\
    #2_{#31}    & #2_{#32}    & \cdots & #2_{#3#4}
  \end{pmatrix}
  \]
}

% 简化版命令（默认矩阵名为A，元素符号为a）: \quickmatrix{行数}{列数}
\newcommand{\quickmatrix}[2]{\flexmatrix{A}{a}{#1}{#2}}

\begin{document}
\title{注释 14.1}
\author{张志聪}
\maketitle

% \begin{zremark}
%   使用定理14.3证明例4
% \end{zremark}
% \textbf{证明：}
% 提示，把缺项补上，即系数$a_i = 0, i = 1, 3, 5, \cdots$。

\begin{zremark}
  定理14.4 的证明过程中，指出$\overline{x}$处幂级数绝对收敛的原因。
\end{zremark}
\textbf{证明：}

由实数稠密性可知，存在$x' \in (-R, R)$且$x' > \overline{x}$，
即$|\overline{x}| < |x'|$，
于是利用定理14.1（阿贝尔定理）可得结论。

\begin{zremark}
  疑问：
  定理14.5 的证明过程中，指出$\sum a_n R^n$收敛，
  但阿贝尔定理明明要求函数项级数一致收敛。
\end{zremark}
\textbf{证明：}
这里的$R$是常量，故$\sum a_n R^n$是数项级数（也可以看做特殊的函数项级数，每项都是常数函数），
收敛的数项级数，满足函数项级数对一致收敛的定义。

\begin{zremark}
  定理14.7 只证明了单向；
\end{zremark}
\textbf{证明：}
todo

\begin{zremark}
  证明

  幂函数(2)的和函数为奇（偶）函数，则(2)式不出现偶（奇）次幂的项。
\end{zremark}
\textbf{证明：}

设幂函数(2)的和函数$f(x)$为奇函数。于是，我们有
\begin{align*}
  f(-x) = f(x)
\end{align*}
即
\begin{align*}
  \sum a_n (-x)^n = - \sum a_n x^n \\
  \sum (-1)^n a_n x^n = \sum - a_n x^n
\end{align*}
由定理14.9可知
\begin{align*}
  (-1)^n a_n = - a_n
\end{align*}
当$n$是偶数是，我们有
\begin{align*}
  a_n   & = - a_n \\
  2 a_n & = 0     \\
  a_n   & = 0
\end{align*}
命题得证。

\begin{zremark}
  证明：例7中
  \begin{align*}
    \sum\limits_{n = 1}^{\infty} (-1)^{n - 1} x^n = \frac{x}{1 + x} \, \, , x \in (-1, 1)
  \end{align*}
\end{zremark}
\textbf{证明：}
\begin{align*}
  \sum\limits_{n = 1}^{\infty} (-1)^{n - 1} x^n
   & = x \sum\limits_{n = 1}^{\infty} (-1)^{n - 1} x^{n - 1} \\
   & = x \sum\limits_{n = 1}^{\infty} (-x)^{n - 1}
\end{align*}
右侧求和是等比数列求和，故
\begin{align*}
  x \sum\limits_{n = 1}^{\infty} (-x)^{n - 1}
   & = x \lim\limits_{n \to \infty} \frac{1(1 - (-x)^{n - 1})}{1 - (-x)} \\
   & = x \cdot \frac{1}{1 + x} \\
   & = \frac{x}{1 + x} 
\end{align*}



\end{document}